\section{Verificación de masas participativa}

Del modelo se obtuvieron los períodos y las masas participativas para cada modo, los cuales se presentan en la Tabla \ref{tab:modal}. Se observa que el primer modo tiene un período de 1.259 segundos y una participación de masa del 1.3\% en la dirección UX, aunque este es un modo torsional con una participación significativa del 44.2\% en la dirección RZ, mientras que el segundo modo tiene un período de 1.141 segundos y una participación de masa del 59.3\% en la dirección UX. El tercer modo tiene un período de 0.854 segundos y una participación de masa del 62.1\% en la dirección UY.

\begin{table}[H]
\centering
\caption{Períodos y participación de masas modal}
\label{tab:modal}
\begin{tabular}{ccccccc}
\hline
\textbf{Modo} & \textbf{T [s]} & \textbf{UX} & \textbf{UY} & \textbf{RZ} & \textbf{$\Sigma$UX} & \textbf{$\Sigma$UY} \\
\hline
1  & 1.259 & 0.013  & 0.000  & 0.442 & 0.013 & 0.000 \\
2  & 1.141 & 0.593  & 0.002  & 0.046 & 0.606 & 0.002 \\
3  & 0.854 & 0.002  & 0.621  & 0.000 & 0.608 & 0.623 \\
4  & 0.328 & 0.001  & 0.000  & 0.109 & 0.609 & 0.623 \\
5  & 0.267 & 0.161  & 0.000  & 0.001 & 0.769 & 0.623 \\
6  & 0.209 & 0.000  & 0.159  & 0.000 & 0.770 & 0.782 \\
7  & 0.152 & 0.000  & 0.000  & 0.044 & 0.770 & 0.782 \\
8  & 0.118 & 0.063  & 0.000  & 0.000 & 0.833 & 0.782 \\
9  & 0.098 & 0.000  & 0.061  & 0.000 & 0.833 & 0.843 \\
10 & 0.093 & 0.000  & 0.000  & 0.027 & 0.833 & 0.843 \\
11 & 0.082 & 0.002  & 0.000  & 0.004 & 0.835 & 0.843 \\
12 & 0.071 & 0.038  & 0.000  & 0.000 & 0.872 & 0.843 \\
13 & 0.065 & 0.000  & 0.000  & 0.018 & 0.873 & 0.844 \\
14 & 0.062 & 0.000  & 0.037  & 0.000 & 0.873 & 0.880 \\
15 & 0.049 & 0.024  & 0.000  & 0.008 & 0.896 & 0.880 \\
\hline
16 & 0.049 & 0.005  & 0.000  & 0.008 & \textbf{0.901} & 0.880 \\
\hline
17 & 0.044 & 0.000  & 0.024  & 0.000 & 0.901 & \textbf{0.904} \\
\hline
18 & 0.040 & 0.010  & 0.000  & 0.036 & 0.911 & 0.904 \\
19 & 0.038 & 0.028  & 0.000  & 0.010 & 0.939 & 0.904 \\
20 & 0.035 & 0.020  & 0.000  & 0.118 & 0.959 & 0.904 \\
\hline
\end{tabular}
\end{table}

En los modos 16 y 17, se alcanza una participación de masa acumulada del 90.1\% en la dirección UX y del 90.4\% en la dirección UY, respectivamente. Esto indica que el modelo es adecuado para representar el comportamiento sísmico del edificio, ya que se cumple con el criterio de participación de masa mínima segun la NCH433. Además, los períodos obtenidos son razonables para un edificio de esta altura y tipo de construcción.